% Pursuing Stacks (cotes 134-2 to 134-8): archive cote 134-3, site page 7 (PDF page 8), typescript page 7.
% typescript. Draft compared: ps2.tex:281-343.
% AI-transcribed by Claude Opus 5.5 (claude-opus-5-5, adaptive thinking, effort high), Message Batches API, batch ps1 (msgbatch_01Rv5hCaNTtUt4JHx1yqoHPL), 2026-11-12; not yet checked by eye.
\scanpage{134-3/7}
\mspage{7}

\margin{7}

\margin{16'.} 16.\hand{\deleted{\unsure{bis}}} One last comment before diving into more technical matters. Without
even climbing up the ladder of increasing sophistication, leading up
to $\infty$-stacks, there is on the very first step, namely with just usual
categories, the possibility of describing arbitrary homotopy types.
Namely, there are two natural, well known ways to associate to a cate-%
gory C (I mean here a ``small'' category, belonging to \deleted{a}\inserted{the} given universe
we are working in) some kind of topological object, and hence a homotopy
type. One is by associating to C the topos $\overset{\hand{$\wedge$}}{C}$ or Top(C), the sheaves on
which are just the presheaves on C, namely arbitrary contravariant func-%
tors from C to (Sets). The other is through the nerve functor\inserted{(7)}, associa-%
ting to C a semi-simplicial set - and hence, if this suits us better,
a topological space, by taking the geometrical realization. \over{b}{B}y a con-%
struction of Verdier, any topos and therefore Top(C) gives rise canoni-%
cally to a pro-object in the category of semi-simplicial sets, and hence
by ``localization'' to a pro-object in the homotopy category (namely a
\hand{``}pro-homotopy type\hand{''} in the terminology of Artin-Mazur). In the same way,
the nerve $\underline{N}$(C) gives rise to a homotopy type - and of course the two
(pro)-homotopy types thus obtained coincide essentially, and may be
called \uline{the} homotopy type of C. When C is a groupoid, we get merely
a 1-truncated homotopy type, namely with homotopy groups \hand{$\pi$}$_n$ which va-%
nish for n\hand{$\geqslant$}2, or equivalently, with connected components (correspon\xout{d}d-%
ding of course to connected components of C) K(\hand{$\pi$},1) spaces. This had
led me \inserted{at one moment} in the late sixties to \inserted{hastily} surmise that even for arbitrary C, we got
merely such sums of K(\hand{$\pi$},1) spaces (namely, that the homotopy type of C
does not change when replacing C by the universal enveloping groupoid,
deduced from C by \xout{xxxxx} making formally invertible all it's arrows).
As Quillen pointed out to me, this is definitely not so - indeed,
using categories, we get (up to isomorphism) arbitrary homotopy types.
This is achieved, I guess, using the \inserted[typed]{left} adjoint functor \inserted{N'} from the
inclusion functor
\[ \diagram{\begin{tikzcd} \text{(Cat)} \ar[r, hook, "\underline{N}"] & \text{(Ss sets)} \end{tikzcd}} \]
which is fully faithful \hand{(}the adjoint functor being therefore a localiza-%
tion functor\hand{)}, and showing that for a semisimplicial set K, the natural
\margin{\unsure{5.3.} This is false, see \unsure{§24} below (p. 21-22) (!8)}
map
\[ \deleted{\illegible{1}}\ K \hand{$\longrightarrow NN'(K)$} \]
is a weak equivalence\deleted{\illegible{1}}\hand{\unsure{;|,}} or what amounts to the same, that the set of
arrows in (Ss sets) by which we localize in order to get (Cat) \hand{(}namely
those transformed into invertible arrows by $\underline{N}$') \deleted{are} is made up with
weak equivalences only\inserted{(!8)}. This would imply that we may reconstruct the
usual homotopy category, up to equivalence, in terms of (Cat), by
